\documentclass[11pt]{article}
\usepackage{amsmath,amssymb,amsthm}
\usepackage{hyperref}

\title{Literature Notes: Operators on Cylindric Partitions}
\author{Notes}
\date{}

\begin{document}
\maketitle

\section{Korff--Palazzo: $h$-positivity for cylindric reverse plane partitions}

\textbf{Reference:} C.~Korff and D.~Palazzo,
\emph{Cylindric symmetric functions and positivity},
\href{https://arxiv.org/abs/1804.05647}{arXiv:1804.05647} (2020).
Published in \emph{Algebraic Combinatorics} 3(1), 191--247.
(Conference version: \href{https://arxiv.org/abs/1802.08977}{arXiv:1802.08977}, FPSAC 2018.)

\medskip
\noindent\textbf{Their setup:}
\begin{itemize}
  \item Objects: cylindric \emph{reverse} plane partitions (CRPPs),
    entries weakly \emph{increasing} along rows and down columns,
    with cylindric periodicity.  Parameters: $k$~rows (our rank),
    $n = $ rank of $\mathfrak{sl}_n$ (bounds part sizes),
    $d = $ winding number.
  \item Generating function: the \emph{cylindric complete symmetric
    function} $h_{\lambda/d/\mu}(x_1, x_2, \ldots)$, a symmetric
    function in infinitely many variables.
\end{itemize}

\medskip
\noindent\textbf{Their main result (Theorem~1):}
The cylindric complete symmetric function expands in the
$\{h_\nu\}$ basis (ordinary complete homogeneous symmetric functions)
with non-negative integer coefficients:
\[
  h_{\lambda/d/\mu} = \sum_\nu N^{\lambda}_{\mu,\nu}\, h_\nu,
  \qquad N^{\lambda}_{\mu,\nu} \in \mathbb{Z}_{\geq 0}.
\]
The coefficients $N^{\lambda}_{\mu,\nu}$ count pairs of permutations
in the wreath product $G(n,1,k) = C_n \wr S_k$.

\medskip
\noindent\textbf{Proof technique:} Entirely representation-theoretic.
Extended affine symmetric group, Verlinde algebra, modular $S$-matrix.
No combinatorial operators are defined.

\medskip
\noindent\textbf{Relationship to our work:}
\begin{itemize}
  \item \textbf{Different ``$h$'':}
    Their $h$-positivity is a symmetric function expansion statement
    ($h_{\lambda/d/\mu}$ expands nonneg in the $\{h_\nu\}$ basis).
    Our $h_m = (q)_m \cdot g_m \geq 0$ is about a single-variable
    polynomial.  These are \textbf{not the same} statement.
  \item \textbf{Wrong direction:}
    Even if their $h$-positivity implies our $h_m \geq 0$ via
    principal specialisation, this is the \textbf{easy} direction.
    The corrected $q$-binomial inverse
    $h_m = \sum_j q^j \binom{m}{j}_q Q_{m-j}$ already shows
    $Q_n \geq 0 \Rightarrow h_m \geq 0$ trivially.  What we need
    is $Q_n \geq 0$, and $h_m \geq 0$ does not imply this
    (the forward formula is an alternating sum).
  \item \textbf{Reverse vs.\ standard:}
    They work with CRPPs (entries non-decreasing), we work with
    CPPs (entries non-increasing).  Related by
    complementation/conjugation but not identical.
  \item Their $k = 3$ matches our rank~3.
  \item Their $n$ is related to our max entry bound, but the
    precise dictionary needs work.
\end{itemize}

\medskip
\noindent\textbf{TODO:}
\begin{enumerate}
  \item Read the full 63-page journal version (arXiv:1804.05647).
  \item Determine whether the principal specialisation of their
    Theorem~1 implies our $h_m(q) \geq 0$.
  \item Check whether their $N^{\lambda}_{\mu,\nu}$ for $k = 3$
    relate to our $Q_{n,c}(q)$ coefficients.
\end{enumerate}


\section{Tingley: crystal operators on cylindric partitions}

\textbf{Reference:} P.~Tingley,
\emph{Three combinatorial models for $\widehat{\mathfrak{sl}}(n)$
crystals, with applications to cylindric plane partitions},
\href{https://arxiv.org/abs/math/0702062}{arXiv:math/0702062} (2007).

\medskip
\noindent
Tingley defines Kashiwara crystal operators $e_i, f_i$
($i = 0, 1, \ldots, n-1$) directly on cylindric plane partitions.
These operators:
\begin{itemize}
  \item \textbf{Preserve} the cylindric structure (profile/interlacing),
  \item \textbf{Change weight} by a simple root $\alpha_i$ (add/remove
    a single box),
  \item Are \textbf{not} weight-preserving.
\end{itemize}

\noindent\textbf{Relevance:} These are the only known local operators
on CPPs that preserve the profile.  They move a single box (unlike
crush/peel which move an entire ring).  Understanding how they
maintain interlacing may be key.


\section{Other relevant references}

\begin{itemize}
  \item \textbf{Corteel--Dousse--Uncu} (2022),
    \href{https://arxiv.org/abs/2201.03047}{arXiv:2201.03047}:
    Layer-peeling gives coupled $q$-difference equations between
    GFs for different profiles.  Algebraic, not bijective.
  \item \textbf{Garsia--Milne} (1981): Involution principle for
    constructing weight-preserving bijections from sign-reversing
    involutions.  Universal but opaque.  Applicable to the $Q_n$
    alternating sum.
  \item \textbf{Warnaar} (2023),
    \href{https://arxiv.org/abs/2111.07550}{arXiv:2111.07550}:
    Proves the $A_2$ Andrews--Gordon identities using CPPs.
    Conjecture~2.7 ($Q_{n,c} \geq 0$) remains open.
\end{itemize}


\end{document}
